\section{Information recovered from public play}
\label{sec:public-information}
The fixed-observer critic uses 561 viewer-relative features
(Appendix~\ref{sec:observer-selfplay}). Public history supports two concrete
reconstructions. A viewer-seat bit makes
the critic representation distinguish actor inputs that it previously merged.
Completed hand discards reveal enough information to reconstruct an opponent's
earlier inputs and identify it within a fixed model class. We verify these
information properties before testing their use in learning.

\subsection{Aligning actor and critic inputs}
The actor and critic also differ in their input representations. The published
actor records the absolute dealer seat, whereas the 561-input critic records
dealer status relative to its viewer. Exchanging seat labels preserves the
critic input but changes two actor feature rows. On 256 retained source
decisions per parent, this relabeling changes 4 and 5 greedy actions,
respectively. All nine changes occur at discard decisions. Legal actions and
public endgame assistance are unchanged, and the actor features agree with
the pinned upstream implementation on all 1,024 original/relabeled probes.

A public viewer-seat bit resolves this ambiguity. With that bit and the
relative-dealer feature, an observation-only decoder recovers the absolute
dealer and exactly reconstructs the complete $300\times50$ actor input from
the 562-element critic row on all 1,024 probes. The probes span discard,
hand-capture, draw-capture and stop/continue decisions. Reconstruction uses
only that active viewer's features, without stepping the saved states or
running the actor. This diagnoses a concrete actor--critic information
mismatch and supplies an explicit repair.

\subsection{Retaining earlier decision information}
\label{sec:temporal-information}
We next test whether later critic features retain the input and action of an
earlier decision. An observation-only decoder receives a later 562-element
row and the public round, turn and phase of
an earlier decision by the same viewer, and reconstructs that decision's actor
input and chosen action. The viewer's initial hand is recovered from its
remaining cards and recorded hand plays. Starting from the recorded initial
field, the decoder replays public card, capture and continuation events up to
the requested address. It receives no simulator state, hidden opponent hand,
stock order, earlier observation or policy output.

All 3,764 chronological pairs within the same played round in the saved
trajectories match exactly: 1,864 pairs for the first parent and 1,900 for the
second. In 1,968 pairs the later actor is the opponent. The checks cover all
four earlier decision phases and separations up to 19 decisions. They reuse
the 512 source positions and add no simulator steps or policy inference.

For a fixed differentiable policy, the score
$\nabla_\theta\log\pi_\theta(a\mid o)$ is determined by its input and chosen
action. Recovering both arguments at a supplied public address determines the
earlier score on these checked pairs. This is measured within-round retention, rather than a
proof for every trajectory or across round boundaries. Stop actions end a round
and have no same-round successor. Appendix~\ref{sec:seat-comparison} tests whether this representation repair
improves playing strength.

\begin{samepage}
\subsection{Recovering opponent information after a hand is revealed}
\label{sec:public-revelation}
A player's eight hand cards are discarded one by one; stock draws never enter
the hand. Once all eight hand discards have been observed, their set reveals
that round's initial hand. Combining it with the public board, captures and
action history reconstructs the player's earlier inputs. This reconstruction
becomes available after the eighth discard, including when that action ends
the round. It cannot be used to choose any earlier action.
\end{samepage}

We retain the two fixed assisted CPU opponents as candidate models. At each
revelation, we reconstruct their inputs at earlier opponent decisions and
remove any candidate whose predicted action disagrees with the observed one.
The candidate set carries across rounds. An unresolved set remains unresolved;
an empty set records model mismatch. The true opponent label is used only for
scoring the diagnostic, never for prediction.

Across all 2,048 retained training prefixes, the procedure identifies the true
opponent in 1,746 cases and leaves 302 unresolved. It makes no incorrect
identification or exclusion of the true model. At the 1,856 natural final-round
starts, 1,598 have an identified opponent (86.1\%;
Table~\ref{tab:public-revelation}). All 36,977 reconstructed historical inputs
and assistance masks agree with the original upstream builders. Independent
replay verifies the time of each revelation and every candidate exclusion.
These prefixes share 512 deals; they are not 2,048 independent full matches.

\begin{table}[ht]
\centering\small
\begin{tabular}{llrrr}
\toprule
Parent & Opponent & Identified & Unresolved & Identified at final root \\
\midrule
1 & WP & 428/512 & 84 & 410/483 \\
1 & Point & 435/512 & 77 & 387/455 \\
2 & WP & 438/512 & 74 & 402/460 \\
2 & Point & 445/512 & 67 & 399/458 \\
\bottomrule
\end{tabular}

\caption{Identification from public training histories within the fixed
WP/Point candidate class. Every identified case is correct. The last column
conditions on reaching the final round; all earlier terminations remain in
the preceding columns.}
\label{tab:public-revelation}
\end{table}

Identification assumes that the opponent belongs to the specified deterministic
class. Unlisted or changing policies can violate that assumption.

\subsection{Playing-strength tests of the recovered information}
The seat reconstruction and opponent identification are exact on the checked
inputs. We test their value for learning with public and masked information
conditions that share architecture, initialization and training schedules.
Every comparison evaluates complete matches from their natural beginnings;
conditional training tails are counted separately.

Adding the public-seat bit yields one positive lower win-change bound among
eight primary comparisons in 160,000 evaluation matches. That gain is against
Point at parent 2. All four public-minus-masked win intervals include zero,
so the comparison does not attribute a replicated gain to the added bit.
Appendix~\ref{sec:seat-comparison} reports the controls and secondary outcomes.

For opponent memory, three studies test progressively earlier and more varied
adaptation: the final round, all rounds after identification, and the same
public starting contexts with freshly sampled cards and future chance.
They evaluate 24,576, 24,576 and 40,960 full matches, respectively, on separate
shared deal sets. Every primary win interval includes zero in all three
studies. Earlier activation and renewed chance therefore do not produce a
replicated improvement under these schedules. Appendix~\ref{app:memory-learning}
retains every primary comparison and the secondary point findings.

The distinction is operational: public reconstruction supplies information
that the original representation or memory omitted; the tested learners have
yet to turn that information into a consistent additional match-win gain.
